% Direct-authority diplomatic French for printed p.213.
\clearpage
\SourcePage{213}

\LayerTarget{tohoku:proof:lemma-5.5.5}
\noindent $U_{s,\mathbf g}\cap f^{-1}(W^s_y)$ soit nulle : un tel $W^s_y$
existe grâce à l'hypothèse sur $P$, au fait qu'il n'y a qu'un nombre fini de
systèmes $\mathbf g$ à considérer, et que les
$U_{s,\mathbf g}\cap f^{-1}(W)$, quand $W$ parcourt un système fondamental de
voisinages de $y$, parcourt un système fondamental de voisinages de l'ensemble
fini des $x\in U_{s,\mathbf g}$ se projetant sur $y$. $Y$ étant paracompact,
il résulte maintenant d'un lemme bien connu en théorie de Čech, qu'on peut
trouver un recouvrement $\mathcal V'=(V'_j)_{j\in J}$ de $Y$, une application
$\varphi':J\to I$ telle que $V'_j\subset U'_{\varphi'(j)}$ pour tout $j\in J$,
et telle que pour tout $n$-simplexe $t$ de $J$, $V'_t$ soit contenu dans un au
moins des ensembles $W_y^{\varphi'(t)}$. Soit alors
$V_j=U_{\varphi'(j)}\cap f^{-1}(V'_j)$, soit
$\mathcal V=(g\cdot V_j)_{(g,j)\in G\times J}$ et considérons l'application
$\varphi:G\times J\to G\times I$ définie par $\varphi'$. On voit aussitôt que
$(\mathcal V,\varphi)$ satisfait aux conditions voulues, et notamment
$\varphi(F)=0$ d'où $\overline\varphi(f^n)=0$.

Le faisceau associé au préfaisceau $H^q(A)$, $q>0$, est nul (lemme 3.8.2). Le
lemme 5.5.5 et la suite spectrale de terme initial 5.5.15) montrent alors :

\LayerTarget{tohoku:theorem:5.5.6}
\noindent\textsc{Théorème 5.5.6.}\quad \textit{Si $G$ est un groupe discret
d'homéomorphismes de $X$, et si $Y=X/G$ est paracompact, alors les
homomorphismes}

\LayerTarget{tohoku:equation:5.5.14}
\[
\tag{5.5.14}
\check H^n(X;G,A)\longrightarrow H^n(X;G,A)
\]

\noindent\textit{sont des isomorphismes pour tout $G$-faisceau abélien $A$.}

\LayerTarget{tohoku:remark:2-low-degree-cech-consequences}
\noindent\textsc{Remarques II.}\quad De la suite spectrale de terme initial
(5.5.15) on tire, sans aucune hypothèse sur $X,G,A$, que

\LayerTarget{tohoku:equation:cech-h1-identification}
\[
\check H^1(X;G,A)=H^1(X,G,A)
\]

\LayerTarget{tohoku:comparison:cech-h2-injection}
\noindent (et que $\check H^2(X;G,A)\to H^2(X;G,A)$ est injectif, car on
vérifie que $E_2^{0,n}=0$ pour tout $n>0$).

\LayerTarget{tohoku:interpretation:g-a-bundles}
Or $\check H^1(X;G,A)$ peut s'interpréter géométriquement, comme on constate
facilement, comme l'ensemble des classes de $G$-$A$-fibrés sur $X$ (i. e. de
fibrés sur $X$ ``à faisceau structural $A$'' [11], et où $G$ opère de façon
compatible avec cette structure). D'ailleurs, $\check H^1(X;G,A)$ est défini
aussi, et admet la même interprétation géométrique, si $A$ est un $G$-faisceau
de groupes non nécessairement commutatifs. Les développement de [11], en
particulier du Chap. V, peuvent alors se transporter à peu près mot à mot dans
le contexte plus général des $G$-fibrés. Il serait d'ailleurs indiqué de donner
un traitement d'``Algèbre Homologique non commutative,'' dans le contexte des
foncteurs et catégories, qui englobe à la fois cette théorie des fibrés, et le
mécanisme algébrique des questions d'extensions de groupes (de Lie et autres)
tel qu'il se trouve développé par exemple dans des papiers de Hochschild et de
A. Shapiro [12, 13, 19].

\LayerTarget{tohoku:section:5.6-equivariant-ext}
\noindent\textbf{5.6.\quad Les $\operatorname{Ext}^n_{O,G}(X;A,B)$.}
Dans ce numéro, $O$ désigne un $G$-faisceau d'anneaux fixé, et on considère la
catégorie abélienne $C^{O(G)}$ des $G$-$O$-Modules (cf. 5.1).

\LayerTarget{tohoku:definition:equivariant-o-homomorphisms}
Le groupe des $G$-$O$-homomorphismes d'un $G$-$O$-Module $A$ dans un autre
$B$ sera noté $\operatorname{Hom}_{O,G}(A,B)$, ou
$\operatorname{Hom}_{O,G}(X;A,B)$ si on veut expliciter

\TohokuContinuation{214}
